% Generated draft: requires mathematical and source review.
\documentclass[11pt]{article}
\usepackage[margin=25mm]{geometry}
\usepackage{fontspec}
\setmainfont{Latin Modern Roman}
\usepackage{amsmath,amssymb,mathtools,unicode-math}
\setmathfont{Latin Modern Math}
\usepackage{enumitem,needspace,xurl,hyperref}
\setlength{\parindent}{0pt}
\setlength{\parskip}{5pt}
\setlength{\emergencystretch}{3em}
\newcommand{\N}{\mathbb{N}}
\newcommand{\Z}{\mathbb{Z}}
\newcommand{\Q}{\mathbb{Q}}
\newcommand{\R}{\mathbb{R}}
\newcommand{\C}{\mathbb{C}}
\newcommand{\F}{\mathbb{F}}
\newcommand{\A}{\mathbb{A}}
\newcommand{\PP}{\mathbb{P}}
\newcommand{\E}{\mathbb{E}}
\newcommand{\eps}{\varepsilon}
\newcommand{\dd}{\,\mathrm{d}}
\begin{document}
\section*{UnsolvedMath ID: 20001300}
\textbf{Source problem code:} catalogue entry
\subsection{Definitions and mathematical statement}
Let \(K\subset\R^n\) be an origin-symmetric convex body with nonempty interior. A zonoid is a Hausdorff limit of finite Minkowski sums of line segments. Define
\[
 \lambda(K)=\inf\{\lambda\ge1:\exists\text{ zonoid }Z,
 \ K\subseteq Z\subseteq\lambda K\}.
\]
For \(C_n=\{x:\sum_i|x_i|\le1\}\), the supplied aid recovers the source's split display as
\[
 \lambda(C_n)=\frac n{2^{n-1}}
 \binom{n-1}{\lfloor(n-1)/2\rfloor}
 \sim\sqrt{\frac{2n}{\pi}}.
\]
The original OCR loses the fraction/root layout and appends page number 12. R. Schneider's question is whether, in each dimension,
\[
 \lambda(K)\le\lambda(C_n)\qquad\text{for every such }K.
\]
It asks for maximality, not uniqueness of maximizers.

The supplied aid reports an attained minimax formulation. Let \(h_K(u)=\sup_{x\in K}\langle u,x\rangle\) and, for a finite even positive Borel measure \(\mu\) on \(\mathbb S^{n-1}\), set
\[
 (C\mu)(u)=\int_{\mathbb S^{n-1}}|\langle u,v\rangle|\,\dd\mu(v),
 \quad m_\mu=\min_u\frac{C\mu(u)}{h_K(u)},\quad
 M_\mu=\max_u\frac{C\mu(u)}{h_K(u)}.
\]
Restrict to \(m_\mu>0\). The reported exact formula is
\[
 \lambda(K)=\min_\mu M_\mu/m_\mu,
\]
with an optimizer and, after scaling, \(m_\mu=1\). The measure convention absorbs any usual factor in the cosine transform. It reports that an optimal sandwich zonoid can be centered, that \(\lambda(TK)=\lambda(K)\) for invertible linear \(T\), and that general affine-invariance language is interpreted with recentering at the symmetry center.

For every compact linear symmetry group of \(K\), an optimal zonoid is reported to be choosable invariant by Minkowski averaging. After coordinates are changed so the group is orthogonal, its generating measure may also be chosen invariant. This is a symmetry reduction, not a solution of the global maximum problem.

Additional supplied benchmarks are \(\lambda(K)\le\sqrt n\) and \(\lambda(K)=1\) for \(n\le2\). The Rademacher zonotope with support function
\[
 h_Z(u)=2^{-n}\sum_{\varepsilon\in\{-1,1\}^n}|\langle u,\varepsilon\rangle|
\]
provides the cross-polytope sandwich, while Schneider's sharp sign-sum inequality is reported to give its optimality and the displayed exact value. Full statements and citations are absent. The aid says the cross-polytope extremality question remains open in the literature it checked; that is a supplied assessment, not a fresh current-status claim.
\subsection{Short English statement}
The source asks whether the cross-polytope is hardest to sandwich by a zonoid up to dilation. The supplied aid gives an attained cosine-transform optimization, symmetry reductions, and known benchmarks, while leaving the general maximality question unresolved.
\subsection{Sources}
\begin{itemize}
\item Local UnsolvedMath registry, permanent ID \(20001300\), source code \texttt{catalogue entry}; \texttt{aim-convex-geometry-notes.json}, \texttt{notes[31]}, tag \texttt{problem}.
\item AIM workshop ``Fourier analytic methods in convex geometry,'' item 2, R. Schneider: \url{https://aimath.org/WWN/fourierconvex/fourierconvex.pdf}. Statement status \texttt{exact}; the supplied verification identifies split fraction/binomial/root formatting and the appended page number. No PDF was accessed for this draft.
\item Schneider's sharp sign-sum inequality and the Rademacher-zonotope benchmark are named in the aid without a complete authoritative citation. No original remarks or literature field are supplied.
\item Research attempt 2: \texttt{partial\_result}, problem status at run \texttt{open}, rigor label \texttt{proved}, novelty confidence low; registry status \texttt{partially\_solved}, classification \texttt{PARTIAL-PROGRESS}. Created and updated 14 August 2026. This machine-generated aid is not peer reviewed; independently verify before citation.
\item Ulam AI / UnsolvedMath Contributors; CC BY 4.0.
\end{itemize}
\end{document}
